% All plots use the 930-point frozen Challenge 2 grid. GapClosed is undefined
% in routing-sufficient points and is rendered gray rather than summarized as 0.

\begin{figure*}[t]
  \centering
  \includegraphics[width=0.86\textwidth]{results/challenge2/figure1_routing_gap_heatmap.pdf}
  \caption{Routing Gap across the predeclared CI-ratio and cross-region-capacity grid (31 CI points per panel).}
  \label{fig:routing-gap-heatmap}
\end{figure*}

\begin{figure*}[t]
  \centering
  \includegraphics[width=0.86\textwidth]{results/challenge2/figure2_gap_closed_heatmap.pdf}
  \caption{CBGP GapClosed on the same network-sensitivity grid. Gray cells are routing-sufficient points where GapClosed is intentionally undefined.}
  \label{fig:gap-closed-heatmap}
\end{figure*}

\begin{figure}[t]
  \centering
  \includegraphics[width=0.48\textwidth]{results/challenge2/figure3_gap_vs_reduction.pdf}
  \caption{Routing Gap versus achieved CBGP reduction for all 930 synthetic points. The dashed line denotes Oracle-equivalent recovery.}
  \label{fig:gap-vs-reduction}
\end{figure}

\begin{figure}[t]
  \centering
  \includegraphics[width=0.48\textwidth]{results/challenge2/figure4_runtime_comparison.pdf}
  \caption{Per-point wall-clock runtime of deterministic CBGP and the Joint MILP Oracle (930 measurements; log scales).}
  \label{fig:runtime-comparison}
\end{figure}

\begin{figure*}[t]
  \centering
  \includegraphics[width=0.86\textwidth]{results/challenge2/figure5_normalized_carbon.pdf}
  \caption{Scenario-level mean operational compute carbon, normalized pointwise to the routing-only baseline before aggregation.}
  \label{fig:normalized-carbon}
\end{figure*}
